\section{Introduction}
\label{sec:intro}

LLM-as-a-judge has become a dominant paradigm for evaluating generated text~\citep{zheng2023judging}.
Rather than recruiting human annotators, practitioners route evaluation through models like GPT-4, which can approximate human preferences at a fraction of the cost.
But what if these automated judges share a blind spot---systematically discounting text that is genuinely interesting to human readers?

Prior work gives reason to worry.
LLMs exhibit self-preference bias, rating their own outputs higher than competitors'~\citep{panickssery2024llm}.
They prefer lower-perplexity text, favoring the familiar over the novel~\citep{wataoka2024self}.
And LLM outputs converge on similar narrative solutions even across model families~\citep{echoes2024}.
Together, these findings suggest a plausible scenario: LLM judges form a consensus that penalizes deviation from a shared ``boring'' distribution, and the texts they collectively reject---\outsiders---may be exactly the ones humans find most surprising and engaging.

We call this the \textbf{artificial outsider hypothesis}: LLM-written texts that other LLMs rate as uninteresting are more likely to be interesting to humans than texts all LLMs prefer.
If true, this would have serious implications for content curation, recommendation systems, and creative AI evaluation pipelines that rely on automated judges.

We test this hypothesis directly.
Using the \hanna dataset~\citep{chhun2022hanna}---1,056 stories from 11 generators with ground-truth human ratings on surprise and engagement---we collect 3,168 interestingness judgments from three OpenAI models (\gptfo, \gptfomini, \gptfourmini).
We then measure whether LLM consensus inversely predicts human interestingness.

Our results are clear: the hypothesis is \textbf{refuted}.
LLM consensus positively correlates with human interestingness ($\rho = +0.48$, $p < 10^{-62}$), and stories in the top LLM consensus quartile receive 47\% higher human interest ratings than those in the bottom quartile.
A complementary analysis of 200 Chatbot Arena battles~\citep{lmsys2023arena} confirms this alignment: humans chose the LLM-preferred response 77\% of the time.
However, when we control for overall text quality, the correlation drops to near zero ($\rho \approx 0.00$), revealing that LLMs are primarily detecting the same quality signal humans detect rather than exhibiting a unique interestingness bias.

We make the following contributions:
\begin{itemize}[leftmargin=*,itemsep=0pt,topsep=0pt]
    \item We propose and rigorously test the artificial outsider hypothesis---that LLM consensus inversely predicts human interestingness---across 1,056 stories and 3,168 LLM evaluations.
    \item We identify a critical quality confound: after controlling for text quality, LLMs show no residual bias for or against interesting content ($\rho \approx 0.00$), demonstrating that raw LLM-human agreement on interestingness is driven entirely by shared quality perception.
    \item We complement our primary analysis with 200 Chatbot Arena battles, confirming that LLM judges align with human preferences at 77--84\% agreement rates.
    \item We provide actionable guidance for practitioners: LLM-as-a-judge is a reasonable proxy for human interestingness judgments, but researchers should control for quality when studying interestingness-specific signals.
\end{itemize}
